% Source-derived chapter 09: Correction
% Original source: universal-p-adic-gross-zagier-formula
\chapter{Correction}
\label{universal-p-adic-gross-zagier-formula-chapter-09}
\source{universal-p-adic-gross-zagier-formula}

\begin{remark}[Source section]
\label{universal-p-adic-gross-zagier-formula-chapter-09-source}
\source{universal-p-adic-gross-zagier-formula}
This chapter records the source section; no Lean declaration has yet been checked for this material.
\notready
\end{remark}

% The source section title was: Correction
Denote by $S_{p, \rm ns}$ the set of $p$-adic places of $F$ that are nonsplit in $E$.
%.\footnote{All notation is as in the original paper \cite{univ},  and all references not accompanied by an external citation  point to it.}
 For a Hecke character $\chi$ of $E$, we consider the following condition:
$$(\star) \qquad  \text{for each  $v\in S_{p, \rm ns}$,  $v$ is inert in $E$ and $\chi_{v}$ is unramified}.$$
After a correction, the  article \cite{nonsplit} only proves the formula of Theorem B for  (ordinary, locally distinguished, potentially crystalline, non-exceptional) representations $\Pi=\pi\ot\chi$ of trivial weight  \emph{satisfying $(\star)$}. As a consequence, the main results of this paper are affected as follows.
\begin{itemize}
\item Theorem A  holds for $\pi_{0} \ot \chi$ under the extra assumption $(\star)$.
\item Theorem B, as well as Theorem $\text{B}^{\text{ord}}$ of \S~7.1.1, hold for   $\Pi=\pi \ot \chi$ under the extra assumption $(\star)$.
\item Theorem C is not affected.
\item Assume that \emph{no $p$-adic place of $F$ is ramified in $E$}. Define a completed homology module $M_{K^{p}}^{\star}$  as in \S~1.3.1, but with the limit being taken only over subgroups $K_{p}$ containing $N_{\G , 0} \ts \prod_{v \in S_{p, \rm ns}} \OO_{E,v}^{\ts}$. Define a space $\cE_{K^{p}}^{\rm ord, \star}$ as in \emph{loc. cit.} with $M_{K^{p}}$ replaced by $M_{K^{p}}^{\star}$. (This is simply the closure in $\cE^{\rm ord, \star}_{K^{p}}$ of the classical points of $\cE^{\rm ord}_{K^{p}}$ that correspond to representations satisfying $(\star)$.)  Define a $\text{Hida}^{\star}$ family for $(\G\ts \H)'$ to be an irreducible component of $\cE_{K^{p}}^{\rm ord, \star}$.

Then Theorem D and Theorem E hold for \emph{$\text{Hida}^{\star}$ families} rather than $\text{Hida}$ families.
\item Theorem F is not affected.
\item Theorem G is not affected, and the related alternative proof of the theorem of Greenberg--Stevens mentioned in Remark 7.3.4 remains valid.
\end{itemize}

The original proofs still apply. For the  main theorems (B and D), the proofs in \S~7.1 need to replace Lemma 7.1.4 with the statement that \emph{in a $\text{Hida}^{\star}$ family, the classical points of trivial weight that satisfy condition \textup{(p-crys)}  are dense}; the last paragraph in the proof of Lemma 7.1.4 applies verbatim to prove that  statement. (The condition (ram) is no longer relevant, since the corrected  version of \cite[Theorem B]{nonsplit} does not require it.)


\begin{enonce*}[remark]{Errata} In  footnote 12, the citation to \cite{nek-heeg}  should  not point  to \S~I.2 (I am grateful to M.H. Nicole for bringing this to my attention), but to Proposition II.2.4 (2). In the third-last line of the proof of Lemma 4.1.1, the symbol `$\cong$' should be replaced by `$\into$'.
\end{enonce*}
